% =================================================
% Справочный лист
% =================================================

\begin{center}
  {\fontsize{25}{30}\selectfont\bfseries\sffamily\color{Primary}
    Логарифм: в какую степень?}

  \vspace{0.3em}

  {\large\sffamily\color{GrayText}
    Понятие логарифма и переход между двумя формами записи}
\end{center}

\vspace{0.45em}

\section{Зачем понадобилось новое число}

Некоторые показательные уравнения решаются подбором:
\[
  2^x=8 \quad\Longrightarrow\quad x=3,
  \qquad
  2^x=\frac14 \quad\Longrightarrow\quad x=-2.
\]

В обоих случаях правую часть удаётся представить в виде степени числа~2.
Но попробуем решить другое уравнение.

\begin{problembox}[Проблемный вопрос]
  Как точно записать корень уравнения
  \[
    2^x=5?
  \]
\end{problembox}

Число 5 находится между соседними степенями двойки:
\[
  2^2=4<5<8=2^3.
\]
Поэтому корень расположен между 2 и 3. График подтверждает существование
единственного решения, но не даёт его точной привычной записи.

\logmotivationfigure

Для этого числа вводят специальное обозначение:
\[
  x_0=\log_2 5.
\]

\begin{conclusionbox}[Смысл обозначения]
  Запись \(\log_2 5\) обозначает такой показатель степени, при котором
  число 2 превращается в число 5:
  \[
    2^{\log_2 5}=5.
  \]
\end{conclusionbox}

\Needspace{13\baselineskip}
\section{Определение логарифма}

\begin{propertybox}[Определение]
  \textbf{Логарифмом положительного числа \(b\) по основанию \(a\)}
  называют показатель степени, в которую нужно возвести число \(a\),
  чтобы получить число \(b\).

  При \(a>0\), \(a\ne1\), \(b>0\):
  \[
    \boxed{\log_a b=c\quad\Longleftrightarrow\quad a^c=b.}
  \]
\end{propertybox}

\logequivalencefigure

\lognotationfigure

Запись
\[
  \log_a b=c
\]
читают так: «логарифм числа \(b\) по основанию \(a\) равен \(c\)».
Она отвечает на один вопрос.

\logquestionfigure

\begin{notebox}[Условия пока нужно помнить как часть записи]
  Для существования вещественного логарифма требуется
  \[
    a>0,\qquad a\ne1,\qquad b>0.
  \]
  Почему возникают именно эти ограничения, будет подробно разобрано в
  следующем комплекте. Сейчас основная цель --- научиться видеть в
  логарифме показатель степени.
\end{notebox}

\Needspace{12\baselineskip}
\section{Как вычислять простые логарифмы}

\begin{methodintro}[Алгоритм]
  Чтобы вычислить \(\log_a b\):
  \begin{enumerate}
    \item задайте вопрос: «В какую степень нужно возвести \(a\), чтобы
      получить \(b\)?»;
    \item представьте число \(b\) в виде степени с основанием \(a\);
    \item показатель найденной степени и есть значение логарифма.
  \end{enumerate}
\end{methodintro}

\subsection{Положительный целый показатель}

\begin{trainingtask}[Пример 1]
  Вычислим \(\log_2 32\).
  \[
    32=2^5
    \quad\Longrightarrow\quad
    \boxed{\log_2 32=5}.
  \]
  Проверка: \(2^5=32\).
\end{trainingtask}

\subsection{Нулевой показатель}

\begin{trainingtask}[Пример 2]
  Вычислим \(\log_7 1\).
  \[
    1=7^0
    \quad\Longrightarrow\quad
    \boxed{\log_7 1=0}.
  \]
\end{trainingtask}

\subsection{Отрицательный показатель}

\begin{trainingtask}[Пример 3]
  Вычислим \(\log_3 \dfrac1{27}\).
  \[
    \frac1{27}=3^{-3}
    \quad\Longrightarrow\quad
    \boxed{\log_3\frac1{27}=-3}.
  \]
\end{trainingtask}

\subsection{Дробный показатель}

\begin{trainingtask}[Пример 4]
  Вычислим \(\log_{25}5\).
  \[
    5=\sqrt{25}=25^{\frac12}
    \quad\Longrightarrow\quad
    \boxed{\log_{25}5=\frac12}.
  \]
\end{trainingtask}

\subsection{Основание меньше единицы}

\begin{trainingtask}[Пример 5]
  Вычислим \(\log_{\frac12}8\).
  \[
    8=2^3=\left(\frac12\right)^{-3}
    \quad\Longrightarrow\quad
    \boxed{\log_{\frac12}8=-3}.
  \]
\end{trainingtask}

\begin{conclusionbox}[Логарифм может быть разным]
  Значение логарифма может быть положительным, равным нулю, отрицательным
  или дробным. Во всех случаях оно остаётся показателем степени.
\end{conclusionbox}

\Needspace{14\baselineskip}
\section{Переход между двумя формами записи}

Одно равенство можно записать двумя способами. При переходе ничего не
вычисляется заново: меняется только язык записи.

\begin{center}
  \renewcommand{\arraystretch}{1.38}
  \begin{tabularx}{\textwidth}{
    >{\centering\arraybackslash}p{0.38\textwidth}
    >{\centering\arraybackslash}p{0.10\textwidth}
    >{\centering\arraybackslash}X}
    \toprule
    \rowcolor{TableHeader}
    Степенная запись & & Логарифмическая запись \\
    \midrule
    \(2^5=32\) & \(\Longleftrightarrow\) & \(\log_2 32=5\) \\
    \(10^{-3}=0{,}001\) & \(\Longleftrightarrow\) & \(\lg 0{,}001=-3\) \\
    \(25^{\frac12}=5\) & \(\Longleftrightarrow\) & \(\log_{25}5=\frac12\) \\
    \(\left(\frac13\right)^{-2}=9\) & \(\Longleftrightarrow\) &
      \(\log_{\frac13}9=-2\) \\
    \bottomrule
  \end{tabularx}
\end{center}

Здесь \(\lg b\) --- сокращённая запись десятичного логарифма:
\[
  \lg b=\log_{10}b.
\]

\begin{methodintro}[Как переводить запись]
  В равенстве \(a^c=b\):
  \begin{itemize}
    \item основание степени \(a\) становится основанием логарифма;
    \item результат возведения \(b\) становится аргументом логарифма;
    \item показатель \(c\) становится значением логарифма.
  \end{itemize}
\end{methodintro}

\begin{warningbox}[Не меняйте роли чисел]
  Из равенства \(2^3=8\) следует
  \[
    \log_2 8=3,
  \]
  а не \(\log_8 2=3\). Основание степени остаётся основанием логарифма.
\end{warningbox}

\begin{trainingtask}[Проверьте себя]
  Перепишите в логарифмической форме:
  \[
    4^3=64,
    \qquad
    5^{-2}=\frac1{25},
    \qquad
    16^{\frac14}=2.
  \]
  Затем каждое равенство переведите обратно.
\end{trainingtask}

\Needspace{13\baselineskip}
\section{Простейшие уравнения}

\subsection{Неизвестно число под знаком логарифма}

\begin{trainingtask}[Пример 6]
  Решим уравнение
  \[
    \log_3 x=4.
  \]
  Перейдём к степенной записи:
  \[
    3^4=x.
  \]
  Поэтому
  \[
    \boxed{x=81}.
  \]
\end{trainingtask}

\begin{trainingtask}[Пример 7]
  Решим уравнение
  \[
    \log_2(x-1)=3.
  \]
  По определению логарифма
  \[
    x-1=2^3=8,
  \]
  откуда
  \[
    \boxed{x=9}.
  \]
\end{trainingtask}

\subsection{Неизвестен показатель степени}

\begin{trainingtask}[Пример 8]
  Уравнение
  \[
    5^x=7
  \]
  не решается подбором целого показателя. По определению логарифма его
  точный корень записывается так:
  \[
    \boxed{x=\log_5 7}.
  \]
\end{trainingtask}

\begin{notebox}[Точная запись ответа]
  Число \(\log_5 7\) --- это уже точный ответ, подобно числам
  \(\sqrt2\) или \(\dfrac13\). Десятичное приближение требуется только
  тогда, когда это отдельно указано в условии.
\end{notebox}

\Needspace{12\baselineskip}
\section{Первые следствия из определения}

Так как \(a^0=1\), то
\[
  \boxed{\log_a1=0}.
\]

Так как \(a^1=a\), то
\[
  \boxed{\log_a a=1}.
\]

Более общо, если число уже записано как степень основания \(a\), то
\[
  \boxed{\log_a a^m=m}.
\]

\begin{propertybox}[Главная идея комплекта]
  Логарифм не нужно воспринимать как отдельный загадочный символ.
  Каждый раз заменяйте вопрос
  \[
    \log_a b=?
  \]
  равносильным вопросом
  \[
    a^{?}=b.
  \]
\end{propertybox}

\Needspace{13\baselineskip}
\section{Типичные ошибки}

\begin{warningbox}[Ошибка 1. Вместо показателя записывают результат]
  Неверно: \(\log_2 8=8\).

  Проверка по определению показывает ошибку: равенство означало бы
  \(2^8=8\), что неверно. Правильно:
  \[
    \log_2 8=3,
    \qquad 2^3=8.
  \]
\end{warningbox}

\begin{warningbox}[Ошибка 2. Меняют основание и аргумент местами]
  Из \(3^4=81\) следует \(\log_3 81=4\), а не \(\log_{81}3=4\).
\end{warningbox}

\begin{warningbox}[Ошибка 3. Считают логарифм произведением]
  Запись \(\log_2 5\) не означает \(\log\cdot2\cdot5\). Это единое
  обозначение числа --- показателя в равенстве \(2^x=5\).
\end{warningbox}

\Needspace{12\baselineskip}
\section{Самостоятельная работа}

\begin{practicebox}[Задание 1. Вычислите]
  \[
    \text{а) }\log_2 16;
    \quad
    \text{б) }\log_5 1;
    \quad
    \text{в) }\log_3\frac19;
    \quad
    \text{г) }\log_9 3;
    \quad
    \text{д) }\log_{\frac14}16.
  \]
\end{practicebox}

\begin{practicebox}[Задание 2. Переведите в другую форму]
  \begin{enumerate}
    \item \(7^2=49\);
    \item \(\log_4 64=3\);
    \item \(10^{-4}=0{,}0001\);
    \item \(\log_{27}3=\dfrac13\).
  \end{enumerate}
\end{practicebox}

\begin{practicebox}[Задание 3. Решите уравнения]
  \[
    \text{а) }\log_2 x=6;
    \qquad
    \text{б) }\log_5(x+2)=2;
    \qquad
    \text{в) }3^x=10.
  \]
\end{practicebox}

\begin{practicebox}[Задание 4. Найдите ошибку]
  Ученик написал:
  \[
    \log_4 64=16,
    \quad\text{потому что}\quad 4\cdot16=64.
  \]
  Объясните, какой вопрос задаёт логарифм, и исправьте решение.
\end{practicebox}

\begin{questionsbox}
  \begin{enumerate}
    \item Как связаны равенства \(\log_a b=c\) и \(a^c=b\)?
    \item Какое число является основанием логарифма?
    \item Может ли значение логарифма быть отрицательным? Приведите пример.
    \item Как точно записать корень уравнения \(2^x=11\)?
  \end{enumerate}
\end{questionsbox}

\begin{notebox}[Краткие ответы]
  Задание 1: \(4;\ 0;\ -2;\ \dfrac12;\ -2\).
  Задание 3: \(64;\ 23;\ \log_3 10\).
  В задании 4 правильный ответ: \(\log_4 64=3\), поскольку \(4^3=64\).
\end{notebox}
